\section{客服 Agent 的第一性问题：成本与质量的长期张力}

Sierra 的切入点不是“再做一个更会聊天的模型”，而是解决企业客服历史难题：高质量服务通常意味着高人力成本；成本优化往往伴随质量下降。传统客服系统通过分层坐席、IVR、工单转派缓解矛盾，但并未从根上消除“成本曲线”和“体验曲线”之间的对抗。

讲者举了多个真实场景：电商退换货、家庭安防设备故障、卫星广播订阅切换等。这些任务有三个共同特征：第一，用户往往在“问题已发生”状态下接触客服；第二，问题需要调用真实业务系统并更新状态；第三，一次交互失败会直接放大用户不满和运营成本。

\begin{importantbox}{核心命题}
customer-facing agent 的价值不在“回答了多少问题”，而在“能否完成可验证的业务动作并减少升级率（escalation）”。这意味着系统评价目标天然偏向 outcome，而不是语言表面流畅度。
\end{importantbox}

\begin{knowledgebox}{为什么“看起来像人”不是首要目标}
在客服场景中，用户关心的是“我的问题现在有没有被解决”。一个礼貌、流畅但不能完成系统操作的 Agent，会比一个表达朴素但动作准确的 Agent 更快失去信任。
\end{knowledgebox}

\begin{warningbox}{常见误区}
如果团队把 KPI 设成“平均对话时长”“回复自然度”或“单轮满意度”，很容易把产品优化成“会聊天但不办事”。当业务进入高峰流量，这类系统会显著增加人工兜底负担。
\end{warningbox}

\subsection{本章小结}

这章建立了后文所有技术选择的约束条件：Agent 必须在真实业务系统中稳定完成任务。凡是不能提高可验证解决率、不能降低升级率、不能改善成本结构的技术优化，都不应被视为主路径。

